% Status: Conditional.
% Missing lemma: for fixed k, Delta, and R, sufficiently large local
% cycle rank forces two vertices with distinct FO^k radius-R pointed types.
% Boundary: finiteness of the local-type space does not imply type
% diversity, and FO^k expressibility of an (m,R)-fan is not established
% when m exceeds the available variable bound.
\section*{Cycle--Orbit Splitting (CLR)}

\begin{lemma}[Finite FO$^k$ local type bound]
For fixed $k,\Delta,R$ there are finitely many FO$^k$ radius-$R$ pointed types in $\Delta$-bounded graphs.
\end{lemma}

\begin{proof}
A path of length at most $R$ is definable by an FO formula with $R$ existential quantifiers.
Using at most $k$ variables, one can quantify a bounded tuple witnessing the anchor $u$ and
$m\le k$ distinct neighbors $z_i$, together with bounded paths from each $z_i$ back to the root.
All constraints are local to radius $R$, so the resulting FO$^k$ formula has quantifier rank $O(R)$.
\end{proof}

% Lean stub reference: URF_Core.lean → CLR.finite_local_types

\begin{lemma}[Cycle rank forces fan]
Large $\mathbb F_2$ cycle rank in $B_{2R}(v)$ forces an $(m,R)$ cycle fan at $v$ for $m$ exceeding the FO$^k$ uniform bound.
\end{lemma}

\begin{proof}
The ball $B_{2R}(v)$ has $O(\Delta^{2R})$ edges.
A cycle basis for $Z_1(B_{2R}(v))$ consists of simple cycles of length at most $2R$.
If the $\mathbb F_2$-dimension exceeds the number of possible edge-disjoint anchor paths from $v$ of length $\le R$,
then by pigeonhole there exists a vertex $u$ on such a path through which many independent cycles pass.
These cycles share the initial segment $(v,u)$ and yield an $(m,R)$-fan with
$m \ge \dim Z_1(B_{2R}(v))/\Delta^{O(R)}$.
\end{proof}

\begin{proof}
Fix $k,\Delta,R$. Up to isomorphism, there are finitely many rooted graphs of radius $R$ and degree at most $\Delta$.
FO$^k$ formulas of quantifier rank $\le R$ can distinguish at most finitely many equivalence classes on this finite set.
Hence the number of FO$^k$ radius-$R$ pointed types is finite.
\end{proof}

\begin{lemma}[Missing cycle-rank-to-type-splitting implication]
For fixed $k,\Delta,R$, sufficiently large
$\dim_{\mathbb F_2} Z_1(B_{2R}(v))$ forces the existence of a vertex
$w$ such that
\[
\tp^{\FO^k}_R(G,v) \neq \tp^{\FO^k}_R(G,w).
\]
\end{lemma}

\noindent
This implication is not proved here. Finiteness of the set of
FO$^k$ radius-$R$ pointed types does not imply that distinct types
occur, and no FO$^k$ formula expressing an $(m,R)$-fan is supplied
when $m$ exceeds the available variable bound.

